\begin{Answer}{2.1}
Với trường hợp $p=2$, ta có thể lấy $x=0, y=z=\omega=1$.

Bây giờ ta xét trường hợp $p>2$. Trước tiên ta xét trường hợp $-1$ là đồng dư bình phương modun  $p$,
khi đó tồn tại một số tự nhiên $a$, $0<a<p-1$ sao cho $a^{2}\equiv-1(modp)$.
Bộ $(x,y,z)=(0,1,a)$. Vì $x^{2}+y^{2}+z^{2}=a^{2}+1$  chia hết cho  $p$  nhưng $1+(p-1)^{2}<p^2$
nên tồn tại $\omega\in\left\{1,2,...,p-1\right\} $ sao cho  $x^{2}+y^{2}+z^{2}- \omega.p=0$ .

Tiếp theo, giả sử $(-1)$  không là đồng dư bình phương modun$p$. Ta phải tìm một số $k$ nào đó để cả $k$ và
 $p-k-1$ đều là đồng dư bình phương. Nếu $\frac{p-1}{2}$ là đồng dư bình phương thì chọn $k$=$\frac{p-1}{2}$.
Nếu ngược lại, thì mỗi đồng dư trong số $\frac{p-1}{2}$ các đồng dư bình phương khác không sẽ rơi vào trong
các cặp $\left\{1, p-2\right\}$, $\left\{2, p-3\right\}$,...,$\left\{\frac{p-3}{2}\right\}$, $\left\{\frac{p+1}{2}\right\}$.
Theo nguyên lý Pigeonhole Principle sẽ có một cặp $(k, p-k-1)$ mà cả hai số $k$ và $(p-k-1)$ đều là đồng dư bình
phương như ta đã định tìm.\\
Vì vậy, ta có thể chọn $x, y\in\left\{0, 1,..., \frac{p-1}{2}\right\}$ sao cho $x^{2}\equiv k (modp)$
và  $y^{2}\equiv p-k-1(modp)$. Cho $z=-1$, ta có $x^2+y^2+z^2$ chia hết cho $p$ và  $x^2+y^2+z^2 <p^2$.
Giá trị  $\omega$ sẽ được xác định như ở trường hợp trước.
\end{Answer}
\begin{Answer}{2.2}
Cho $x=y$, ta được $f(0)=0$.\\
Cho $x=-1, y=0$ ta được $f(1)=-f(-1)$.\\
Cho $x=a, y=1$, sau đó cho $x=a, y=-1$ ta có:\\
\centerline {$f(a^{2}-1)=(a-1)\left[f(a)+f(1)\right]$}\\
\centerline {$f(a^{2}-1)=(a+1)\left[f(a)-f(1)\right]$}\\
Cho các vế phải của các phương trình đó bằng nhau và giải phương trình đối với $f(a)$ ta được $f(a)=f(1).a$ với mọi a.\\
Như vậy, mọi hàm số nào thoả mãn ràng buộc đã cho phải có dạng $f(x)=kx$ với hằng số $k$ nào đó.
Ngược lại, bất kỳ hàm số nào có dạng $f(x)=kx$ với hằng số $k$ nào đó rõ ràng đều thoả mãn yêu cầu bài toán.
\end{Answer}
\begin{Answer}{2.3}
Các góc xét đến đều là các góc định hướng ngoại trừ các trường hợp nói khác đi.\\
Giả sử chúng ta có một tam giác tuỳ ý $XYZ$ với tâm đường tròn nội tiếp là điểm $I$ và tâm đường tròn bàng
tiếp $I_{X}$ đối diện với góc $X$. Suy ra $X, I, I_{X}$ thẳng hàng.
Ta có $\widehat{IYI_{X}}$=$\frac{\pi}{2}$= $\widehat{IZI_{X}}$ vì vậy tứ giác $IYI_{X}Z$ là nội tiếp
được và $\widehat{XI_{X}Y}$=$\widehat{II_{X}Y}$=$\widehat{IZY}$ hay $\widehat{YI_{X}X}$=$\widehat{YZI}$.

\begin{center}
\includegraphics[scale =0.7]{hinh3}
\end{center}

Gọi $I_{1}$, $I_{2}$ lần lượt là tâm đường tròn nội tiếp các tam giác $ABD$ và tam giác $ACD$.\\
Từ giả thiết ta suy ra $P$, $Q$ là các tâm đường tròn bàng tiếp của tam giác $ABD$ đối diện với góc $A$ và góc $D$,
tương tự $R$, $S$ là các tâm đường tròn bàng tiếp của tam giác $ACD$ đối diện với góc $A$ và góc $D$.\\
Áp dụng kết quả của phần trên với $(X, Y, Z, I_X)$ là $(A, D, B, P)$, $(D, A, B, Q)$, $(A, D, C, R)$ và $(D, A, C, S)$
ta được, $\widehat{APD}$=$\widehat{I_{1}BD}$, $\widehat{AQD}$=$\widehat{ABI_{1}}$,
$\widehat{ARD}$=$\widehat{I_{2}CD}$, và $\widehat{ASD}$=$\widehat{ACI_{2}}$.\\
Khi coi các góc sau là không định hướng, ta thấy $\widehat{I_{1}BD}$, $\widehat{ABI_{1}}$,
$\widehat{I_{2}CD}$ và $\widehat{ACI_{2}}$ đều bằng $\frac{\widehat{AQD}}{2}$ $=$ $\frac{\widehat{ACD}}{2}$.\\
Hơn nữa, các góc trên đều cùng một hướng, nên nếu coi chúng là những góc định hướng, chúng sẽ bằng nhau.
Như vậy (trở lại với những góc định hướng) ta có: $\widehat{APD}$ $=$ $\widehat{AQD}$ $=$ $\widehat{ARD}$
$=$ $\widehat{ASD}$ và bốn điểm $P, Q, R, S$ cùng nằm trên cung tròn trương bởi $A, D$.
\end{Answer}
\begin{Answer}{2.4}
Với mỗi số thực $x$, đặt\\
\centerline {$\left\{x\right\}$=$x-\left[x\right]$ $\in [0,1)$}.\\
Ta có\\
 \centerline{$\left[\frac{2k^2}{p}\right]$=$\frac{2k^2}{p}$-$\left\{\frac{2k^2}{p}\right\}$ và
$\left[\frac{k^2}{p}\right]$=$\frac{k^2}{p}$-$\left\{\frac{k^2}{p}\right\}$}\\
Ta được\\
\centerline {$\left[\frac{2k^2}{p}\right]$-$2$ $\left[\frac{k^2}{p}\right]$=$2$ $\left\{\frac{k^2}{p}\right\}$
-$\left\{\frac{2k^2}{p}\right\}$}\\
Nếu ${\left\{x\right\}<\frac{1}{2}}$    thì    ${2\left\{x\right\}-\left\{2x\right\}}$=${\left\{x\right\}-2\left\{x\right\}}$=$0$\\
Nếu $\left\{x\right\}\geq\frac{1}{2}$     thì   ${2\left\{x\right\}-\left\{2x\right\}}$=$2{\left\{x\right\}-(2\left\{x\right\}-1)}$=$1$\\
Như vậy, tổng cần tính trong bài ra sẽ bằng $\alpha$ là số các phần tử $k$ trong $\left[1, p-1\right]$ sao cho
$\left\{\frac{k^2}{p}\right\}$ $\geq$ $\frac{1}{2}$, hay bằng với số đồng dư $k$ khác không mà $k^2$ là đồng dư mô
đun $p$ với một số nào đó trong $\left[\frac{p+1}{2}, p-1\right]$.\\
Vì $p$ là số nguyên tố đồng dư với $1$ mô đun $p$, ta đã biết $-1$ $\equiv$ $d^2$ $(mod p)$, với $d$ là một số nào đó.
Phân chia các đồng dư mô đun $p$ khác không thành $\frac{p-1}{2}$ cặp dạng $\left\{a, da\right\}$ sao cho
 $a^2$ $\equiv$ $-(da)^2$ $(modp)$.\\
Vì vậy, có đúng một đồng dư trong mỗi cặp mà bình phương của nó đồng dư với một số nào đó trong
$\left[\frac{p-1}{2}, p-1\right]$, và có tất cả $\frac{p-1}{2}$ đồng dư như thế.\\
Từ đó suy ra tổng đã cho bằng $\frac{p-1}{2}$.
\end{Answer}
\begin{Answer}{2.5}
Trước tiên ta chứng minh rằng nếu $8$ $\backslash$ $mn$, thì hình chữ nhật  $m$x$n$ có thể  được xếp bởi các hình đã cho.\\
Trường hợp 1: Cả $m$ và $n$ đều là số chẵn. Không mất tính tổng quát ta giả sử rằng $4$ $\backslash$ $m$, $2$ $\backslash$ $n$.
Hai hình đã cho có thể ghép được một hình chữ nhật kích thước $4$x$2$, và $m.n$/$8$ hình chữ nhật như vậy sẽ ghép thành
một hình chữ nhật kích thước $m$x$n$ (gồm $n/2$ hàng và $m/4$ cột).\\
Trường hợp 2: Hoặc $m$ hoặc $n$ lẻ. Không mất tính tổng quát, ta giả sử rằng $m$ là số lẻ. Khi đó $8$ $\backslash$ $n$.
Vì $m>1$ nên $m$ $\geq$ $3$. Ta có thể ghép được một hình kích thước $3$x$8$ như hình vẽ sau:\\

\begin{center}
\includegraphics[scale =0.7]{hinh4}
%\input{bai5bkorea.tex}
\end{center}

Những hình $3$x$8$ như vậy có thể ghép thành hình chữ nhật ($3$x$n$).\\
Nếu $m=3$, ta đã ghép xong. Trong trườg hợp ngược lại,  $m>3$, thì phần còn lại $(m-3)$x$n$ có thể ghép như trong trường
hợp $1$ vì $2$$\backslash$$(m-3)$.\\
Bây giờ ta sẽ chứng minh rằng: nếu hình chữ nhật có kích thước $(m$x$n)$ được ghép bởi các hình trên thì $8$$\backslash$$m.n$.
Vì mỗi một hình $L$  có diện tích là $4$ nên $4$$\backslash$$(m.n)$. Không mất tính tổng quát, giả sử $2$$\backslash$$n$, và
tô $m$ hàng trong hình chữ nhật $m$x$n$ thành hai màu đen trắng cạnh nhau. Mỗi mảnh hình chữ $L$ trong hình chữ nhật được
 ghép sẽ gồm một số lẻ ô đen hình vuông. Vì có tất cả $1$ số chẵn $(n$x$\left[\frac{m}{2}\right])$ ô vuông màu đen, nên hình
chữ nhật được ghép chứa $1$ số chẵn các hình chữ $L$, mà ta đặt số đó là $2k$. Như vậy $m.n=8k$, hay $8$$\backslash$$mn$.
\end{Answer}
\begin{Answer}{2.6}
Đặt vế trái của bất đẳng thức là $S$. Vì $a$$\geq$$b$$\geq$$c$>$0$ và $x$$\geq$$y$$\geq$$z$>$0$ nên ta có
 $bz+cy$$\geq$$by+cz$ suy ra\\
\centerline {$(by+cz)(bz+cy)$ $\leq$ $(by+cz)^2$ $\leq$ $2\left[(by)^2+(cz)^2\right]$ }\\
Đặt $\alpha$$=(ax)^2$; $\beta$$=(by)^2$; $\gamma$$=(cz)^2$, khi đó ta có:\\
\centerline {$\frac{a^2x^2}{(by+cz)(bz+cy)}$ $\geq$ $\frac{a^2x^2}{2\left[(by)^2+(cz)^2\right]}$ $=$ $\frac{\alpha}{{\beta}+
{\gamma}}$}\\
Áp dụng tương tự cho hai bất đẳng thức, ta có\\
\centerline {$S$ $\geq$ $\frac{1}{2}$ $({\frac{\alpha}{{\beta}+{\gamma}}}+{\frac{\beta}{{\gamma}+{\alpha}}}
+{\frac{\gamma}{{\alpha}+{\beta}}})$}\\
Sử dụng bất đẳng thức Cauchy-Schwarz, ta có\\
\centerline {$({\frac{\alpha}{{\beta}+{\gamma}}}+{\frac{\beta}{{\gamma}+{\alpha}}}+{\frac{\gamma}{{\alpha}
+{\beta}}})$$({\alpha}({\beta}+{\gamma})+{\beta}({\alpha}+{\gamma})+{\gamma}({\alpha}+{\beta})) {\geq} ({\alpha}
+{\beta}+{\gamma})^2$ }\\
mà vế phải bằng\\
\centerline  {    $\frac{1}{2}$$\left[({\alpha}-{\beta})^2+({\beta}-{\gamma})^2+({\gamma}-{\alpha})^2+3({\alpha}{\beta}
+{\alpha}{\beta}+{\gamma}{\alpha})\right]$ $\geq$ $\frac{3}{2}$ $(2{\alpha}{\beta}+2{\beta}{\gamma}+2{\gamma}{\alpha})$ }\\
Do đó,\\
\centerline {$S$ $\geq$ $\frac{1}{2}$
$({\frac{\alpha}{{\beta}+{\gamma}}}+{\frac{\beta}{{\gamma}+{\alpha}}}
+{\frac{\gamma}{{\alpha}+{\beta}}})\geq\frac{1}{2}\frac{({\alpha}+{\beta}+{\gamma})^2}{(2{\alpha}{\beta}+2{\beta}{\gamma}+2{\gamma}{\alpha})}\geq\frac{3}{4}$                           }\\
Vậy bài toán đã được chứng minh.
\end{Answer}
